\section{Introduction}\label{sec:intro}

\subsection{Two pipelines that disagree about contexts}

Unified randomization keeps a game completable by reasoning relative to a fixed \emph{reference}: a random contextual sort of a world that is already known to be completable. Every choice it makes (a new base for a head, a new ingredient list for a recipe) is accepted only if every promised pebble can still be backed by strictly earlier pebbles of that sort (\file{randomizations/graph/unified/skeleton/promotion.lua}). If nothing better is found, the reference's own choice (the vanilla base, the vanilla ingredients) is always available as a fallback. Mechanic contexts are hard constraints: rooms and automatability are always kept, isolatability where declared (\file{skeleton/protection.lua}).

Planetary randomization moves \emph{features} between planets: ocean tiles and their fluid, resource placements, lightning, freezing (\file{randomizations/planetary/*.lua}). This changes, on purpose, which contexts mechanics have: lava is pumped on another planet, Vulcanus pumps something else. So planetary randomization cannot keep ``every mechanic keeps all its contexts''. It has its own rules instead (\file{randomizations/planetary/check.lua}, \code{protection.planetary_kept_context}): recipes stay reachable somewhere, planet-locked recipes keep their exact contexts, mechanics keep rooms and automatability, isolatability only for rocket building and electricity, and moved features follow their feature.

The two are joined today by running planetary randomization \emph{first}, as a separate stage in \file{data-final-fixes.lua}. It must leave behind a world that satisfies its rules by itself, using scaffold recipes (variants of recipes that used a planet's old fluid), in-place recipe edits, extra resource patches and climate reverts. Unified then takes that world as ``vanilla'': \code{old_data_raw}, the initial sorts and the final MECHCHECK are all taken after the planetary stage (Figure~\ref{fig:pipeline}, top).

\begin{figure}[t]
\centering
\resizebox{0.98\linewidth}{!}{%
\begin{tikzpicture}[
  box/.style={draw, rounded corners=3pt, align=center, font=\small, minimum height=10mm, inner sep=4pt, fill=white},
  ar/.style={-{Stealth[length=2.4mm]}, thick}]
% current
\node[font=\bfseries\small, anchor=west] at (-0.2,1.3) {Today};
\node[box, fill=cvan!10] (v0) at (0.8,0) {vanilla\\$G_0$};
\node[box, fill=cshift!12, text width=31mm] (pl) at (3.9,0) {planetary stage\\shift + scaffolds,\\edits, patches, reverts};
\node[box, text width=19mm] (w1) at (7.2,0) {$G_1{+}S$\\{\scriptsize (new ``vanilla'')}};
\node[box, text width=24mm] (un) at (10.2,0) {unified\\{\scriptsize first pass, promotion}};
\node[box, text width=23mm, font=\footnotesize] (mc) at (13.5,0) {MECHCHECK\\{\scriptsize vs $G_1{+}S$}};
\draw[ar] (v0) -- (pl); \draw[ar] (pl) -- (w1); \draw[ar] (w1) -- (un); \draw[ar] (un) -- (mc);
\node[lbl, below=1mm of pl] {PLANETCHECK vs $G_0$};
% proposed
\node[font=\bfseries\small, anchor=west] at (-0.2,-2.0) {Proposed};
\node[box, fill=cvan!10] (p0) at (0.8,-3.3) {vanilla\\$G_0$};
\node[box, fill=cshift!12, text width=17mm] (ps) at (3.1,-3.3) {shift $\tau$\\{\scriptsize goals $\tau_*$}};
\node[box, text width=22mm] (pr) at (5.75,-3.3) {reference\\{\scriptsize repaired or}\\{\scriptsize superposed}};
\node[box, text width=22mm] (pu) at (8.65,-3.3) {unified\\{\scriptsize debt-aware}\\{\scriptsize promotion}};
\node[box, fill=crep!12, text width=17mm] (pl2) at (11.35,-3.3) {settlement\\ladder};
\node[box, text width=17mm] (pc) at (13.8,-3.3) {check\\{\scriptsize vs $\tau_*\Gamma_0$}};
\draw[ar] (p0) -- (ps); \draw[ar] (ps) -- (pr); \draw[ar] (pr) -- (pu); \draw[ar] (pu) -- (pl2); \draw[ar] (pl2) -- (pc);
\draw[ar, cdebt, dashed] (pr.south) to[bend right=18] node[lbl, below, text=cdebt] {debts (IOUs)} (pl2.south);
\end{tikzpicture}}
\caption{Top: the current order. The planetary stage has to produce a world that is completable on its own, and unified treats it as vanilla. Bottom: the pipeline this report argues for. The shift transports the goals; unified runs against a \emph{reference} that is repaired or superposed; whatever the reference owes (its debts) is settled after unified, and the result is checked against the transported vanilla goals.}
\label{fig:pipeline}
\end{figure}

\subsection{The question}

The user asked for a theoretical basis for multipass systems in which a pass may shift contexts the way planetary randomization does, and for a side investigation into letting later passes fix ``a small number of contradictions'' that such a shift creates. The motivation is integration: unified should be able to fix some of planetary's contradictions itself, instead of the planetary stage settling everything up front with scaffolds. Two facts stand in the way:
\begin{enumerate}[nosep]
\item planetary randomization changes contexts, while unified's hard constraint is that mechanics keep theirs;
\item unified relies on a completable version of the game, both to rank pebbles and as the fallback that makes it always able to finish.
\end{enumerate}

\subsection{What this report does}

\begin{itemize}[nosep]
\item Section~\ref{sec:model} states the logic as a Horn program over pebbles, so that ``reachable'', ``sort'' and ``backing'' have exact meanings, and shows why the order-independent discovery rule (home contexts) matters for everything that follows.
\item Section~\ref{sec:reference} isolates what promotion actually needs: a \emph{reference}, meaning a completable world with a proof order. Vanilla is one reference among many. Multipass is a relay of references.
\item Section~\ref{sec:shift} defines context shifts, goal transport and contradictions. It proves the dilemma behind the user's question: vanilla is a reference but cannot see moved features, and the shifted world sees them but is not a reference.
\item Section~\ref{sec:constructions} gives three ways out: \emph{scaffolded}, \emph{repaired} and \emph{superposed} references. The last one carries explicit debts that are settled after unified, along a ladder of increasingly visible fixes.
\item Section~\ref{sec:repair} is the side investigation. It measures contradictions on real Space Age seeds (goals, causes, repairs) and gives a staged sort that finds small repair sets. It also finds where unified's current move space does and does not cover the repairs.
\item Section~\ref{sec:pipeline} assembles a concrete pipeline and lists the decisions that are the user's to make.
\end{itemize}

Game facts used here were checked against the installed 2.1 game data and prototype documentation (\file{/Applications/factorio.app/Contents/data}, \file{doc-html/prototype-api.json}, version 2.1.20); where one is stated, its source is given next to it.
